\documentclass[aps,prl,twocolumn,superscriptaddress,showpacs]{revtex4-2}

\usepackage{amsmath,amssymb}
\usepackage{graphicx}
\usepackage{bm}
\usepackage{hyperref}
\usepackage{booktabs}
\usepackage{xcolor}

\newcommand{\PB}[2]{\{#1,\,#2\}}
\newcommand{\dd}{\mathrm{d}}
\newcommand{\RR}{\mathbb{R}}
\newcommand{\NN}{\mathbb{N}}
\newcommand{\calA}{\mathcal{A}}
\newcommand{\calF}{\mathcal{F}}
\newcommand{\calH}{\mathcal{H}}

\begin{document}

\title{The Poisson Algebra of Pairwise Interactions:\\
A Calogero-Moser Integrability Test}

\author{[Authors]}
\affiliation{[Institution]}

\date{\today}

\begin{abstract}
The Poisson algebra generated by the pairwise Hamiltonians of a
three-body system carries a natural filtration whose dimension sequence
$d(k)$ is a structural invariant.  For Newtonian gravity and all other
singular pair potentials tested to date, this sequence takes the
universal value $[3, 6, 17, 116, \ldots]$, while the harmonic
oscillator gives a finite algebra with $d(k) = 15$ for $k \geq 3$.
We test whether this dichotomy reflects integrability by computing
$d(k)$ for the one-dimensional Calogero-Moser system with $1/r^2$
interactions---an exactly integrable model that arises naturally in
Galperin's billiard construction for the digits of $\pi$.  We find
$d(k) = [3, 6, 17, 116]$ for the integrable Calogero-Moser model,
identical to non-integrable gravity, and invariant across all Galperin
superintegrable mass ratios.  The Poisson algebra dimension sequence
therefore detects the \emph{singularity class} of the pair
potential---not its integrability properties---reflecting the
combinatorial structure of the pairwise interaction graph rather than
global dynamical conservation laws.
\end{abstract}

\maketitle


% ==================================================================
\section{Introduction}
\label{sec:intro}
% ==================================================================

A remarkable observation connects two seemingly unrelated problems in
mathematical physics.  Galperin~\cite{Galperin2003} showed that
counting collisions between two balls and a wall---a classical
one-dimensional three-body system---yields the digits of $\pi$ when the
mass ratio takes the form $m/M = 10^{2N}$.  Davis, Geyer, and
Ravishanker~\cite{Davis2017} extended this analysis, demonstrating that
the billiard dynamics map exactly onto the Calogero-Moser model with
inverse-square interactions.  In particular, the squared angular
momentum $L^2 = (M+m)^2(Yw - yW)^2$ is a conserved quantity, rendering
the system Liouville integrable, while at special mass ratios
$m/M = \tan^2(\pi/q)$ for integer $q$, a third independent integral
(the Chevalley invariant $J$) appears, yielding superintegrability.

Independently, recent investigations of the three-body problem have
revealed that the Poisson algebra $\calA$ generated by the three
pairwise Hamiltonians $\{H_{12}, H_{13}, H_{23}\}$ carries a rich
algebraic structure.  Using auxiliary variables $u_{ij} = 1/r_{ij}$ to
maintain polynomial expressions, the iterated Poisson brackets of
these generators can be organized into a filtration
$\calA_0 \subset \calA_1 \subset \calA_2 \subset \cdots$, where
$\calA_k$ is spanned by brackets of depth at most $k$.  The dimension
sequence $d(k) = \dim \calA_k$ is a computable structural invariant.

For Newtonian gravity ($V \propto 1/r$) with $N = 3$ bodies, this
sequence is $d(k) = [3, 6, 17, 116, \ldots]$, and it is invariant
under changes in spatial dimension ($d = 1, 2, 3$), mass ratios, and
coupling constants.  Strikingly, \emph{every} singular pair potential
tested---$1/r^2$, $1/r^3$, $\log r$, Yukawa, and
Coulomb~\cite{paper1,paper3}---produces the identical sequence.  In contrast, the harmonic oscillator
($V \propto r^2$) generates a finite algebra with $d(k) = [3, 6, 13,
15, 15]$, saturating by level 3.

This dichotomy raises a natural question: does the Poisson algebra
dimension sequence distinguish integrable from non-integrable systems?
The Calogero-Moser model in one dimension is the ideal test case---it
is exactly integrable, with known conserved quantities, yet its pair
potential $V \propto 1/r^2$ is singular.  If integrability constrains
the growth of the pairwise Poisson algebra, the one-dimensional
Calogero system should show a dimension sequence different from
non-integrable gravity.  If it does not, the sequence reflects
something else entirely: the singularity structure of the interaction.

In this paper, we report the computation of $d(k)$ for the
one-dimensional Calogero-Moser three-body system through filtration
level $k = 3$.  We find $d(k) = [3, 6, 17, 116]$---identical to every
other singular potential.  The result is invariant across all Galperin
superintegrable mass ratios ($q = 3, 4, 5, 6$) and generic
non-superintegrable mass configurations.


% ==================================================================
\section{Method}
\label{sec:method}
% ==================================================================

\subsection{Pairwise Hamiltonians and Auxiliary Variables}

Consider $N = 3$ particles in $d$ spatial dimensions with positions
$\bm{q}_i$ and momenta $\bm{p}_i$ ($i = 1, 2, 3$), interacting through
a pairwise potential $V(r_{ij})$.  The three pairwise Hamiltonians are
\begin{equation}
  H_{ij} = T_i + T_j + V(r_{ij}),
  \qquad T_i = \frac{|\bm{p}_i|^2}{2m_i},
  \label{eq:Hij}
\end{equation}
where $r_{ij} = |\bm{q}_i - \bm{q}_j|$.

For singular potentials $V \propto r_{ij}^{-n}$, Poisson brackets of
the $H_{ij}$ generate expressions involving inverse powers of
$r_{ij}$ and derivatives thereof.  To maintain polynomial structure,
we introduce auxiliary variables $u_{ij} = 1/r_{ij}$ and augment the
phase space with the chain rule relations
\begin{equation}
  \frac{\partial u_{ij}}{\partial q_i^{\alpha}}
  = -u_{ij}^3 \, (q_i^{\alpha} - q_j^{\alpha}),
  \label{eq:chainrule}
\end{equation}
where $\alpha$ indexes spatial components.  In terms of these variables,
$V_{ij} = -g \, m_i m_j \, u_{ij}^n$ and all Poisson brackets reduce to
polynomial expressions in $(\bm{q}, \bm{p}, \bm{u})$.

\subsection{Filtration and Dimension Sequence}

The \emph{filtration level} of a generator is defined recursively:
$H_{ij}$ has level 0, and $\PB{f}{g}$ has level
$\ell(f) + \ell(g) + 1$.  We define
\begin{equation}
  \calA_k = \mathrm{span}\{g : \ell(g) \leq k\},
  \qquad d(k) = \dim \calA_k.
  \label{eq:filtration}
\end{equation}
The dimension sequence $d(0), d(1), d(2), \ldots$ is a structural
invariant of the algebra.

\subsection{Computational Procedure}

At each level $k$, we compute all Poisson brackets between generators
at levels $k_1$ and $k_2$ with $k_1 + k_2 + 1 = k$, simplify the
resulting expressions via \texttt{sympy.expand}, and collect them as
candidate new generators.  The functional independence of all
generators through level $k$ is determined by constructing an
evaluation matrix
\begin{equation}
  M_{si} = g_i(\bm{z}_s), \qquad
  s = 1, \ldots, N_\text{samples},
  \label{eq:evalmatrix}
\end{equation}
where $\{\bm{z}_s\}$ are random phase-space points drawn uniformly
from a compact domain.  The dimension $d(k) = \mathrm{rank}(M)$ is
determined by singular value decomposition (SVD), with the rank
identified by a gap in the singular value spectrum exceeding
$10^7$ between adjacent values.  This numerical rank determination has
been validated against exact symbolic rank computation
(via echelon form over $\mathbb{Q}$) in all cases where the latter is
tractable.

For the Calogero-Moser potential $V_{ij} = -g\, m_i m_j / r_{ij}^2$,
the Hamiltonian in auxiliary variables is
$H_{ij} = T_i + T_j - g\, m_i m_j \, u_{ij}^2$.  The chain rule
table is unchanged because it depends only on the definition
$u_{ij} = 1/r_{ij}$.

All computations were performed using SymPy for symbolic algebra and
NumPy/SciPy for numerical SVD, with $N_\text{samples} = 500$
phase-space evaluation points.


% ==================================================================
\section{Results}
\label{sec:results}
% ==================================================================

\subsection{The Critical Experiment: 1D Calogero-Moser}

The one-dimensional Calogero-Moser system with $N = 3$ equal-mass
particles and $V_{ij} = -1/r_{ij}^2$ is exactly integrable.  Its
Poisson algebra dimension sequence through level 3 is
\begin{equation}
  d(k) = [3, 6, 17, 116].
  \label{eq:CM_result}
\end{equation}
This is \emph{identical} to the sequence obtained for one-dimensional
Newtonian gravity ($V = -1/r$), which is not integrable.  The SVD gap
ratios exceed $10^7$ at each level, confirming the rank determination
is unambiguous (Table~\ref{tab:gaps}).

\begin{table}[b]
\caption{SVD gap ratios for the 1D Calogero-Moser ($1/r^2$) system
at each filtration level.}
\label{tab:gaps}
\begin{ruledtabular}
\begin{tabular}{cccl}
Level $k$ & $d(k)$ & Gap ratio & $\sigma_{d(k)} / \sigma_{d(k)+1}$ \\
\colrule
0 & 3  & $\infty$       & rank 3 of 3 \\
1 & 6  & $\infty$       & rank 6 of 6 \\
2 & 17 & $2.7 \times 10^{14}$ & $1.4 \times 10^{-1}$ / $5.3 \times 10^{-16}$ \\
3 & 116 & $6.9 \times 10^{7}$ & $3.0 \times 10^{-4}$ / $4.3 \times 10^{-12}$ \\
\end{tabular}
\end{ruledtabular}
\end{table}

\subsection{Complete Potential Comparison}

Table~\ref{tab:comparison} summarizes the dimension sequence across
all tested potential types and spatial dimensions.  Every singular
potential---$1/r$, $1/r^2$, and $1/r^3$---produces the identical
sequence $[3, 6, 17, 116]$ in both one and two spatial dimensions.
Companion computations~\cite{paper1,paper3} for $\log r$, Yukawa
($e^{-\mu r}/r$), and Coulomb potentials in two dimensions also yield
this sequence.  The harmonic oscillator ($V \propto r^2$) is the sole
exception, generating a finite algebra that saturates at $d(3) = 15$.

\begin{table}[t]
\caption{Poisson algebra dimension sequences for various
pair potentials and spatial dimensions.  All $N = 3$, equal masses.}
\label{tab:comparison}
\begin{ruledtabular}
\begin{tabular}{llccccc}
$d$ & Potential & Integ.\ & $d(0)$ & $d(1)$ & $d(2)$ & $d(3)$ \\
\colrule
1 & $1/r$       & No  & 3 & 6 & 17 & 116 \\
1 & $1/r^2$     & Yes & 3 & 6 & 17 & 116 \\
1 & $1/r^3$     & No  & 3 & 6 & 17 & 116 \\
\colrule
2 & $1/r$       & No  & 3 & 6 & 17 & 116 \\
2 & $1/r^2$     & No$^\dagger$  & 3 & 6 & 17 & 116 \\
2 & $\log r$$^\S$    & No  & 3 & 6 & 17 & 116 \\
2 & Yukawa$^\S$      & No  & 3 & 6 & 17 & 116 \\
\colrule
2 & $r^2$       & Yes & 3 & 6 & 13 & 15  \\
\end{tabular}
\end{ruledtabular}
\footnotetext{$^\dagger$The $1/r^2$ potential is integrable in 1D
(Calogero-Moser) but not in 2D for scalar particles.}
\footnotetext{$^\S$From companion computations~\cite{paper1,paper3};
not recomputed in this work.}
\end{table}


\subsection{Mass Ratio Invariance at Superintegrable Points}

Galperin's billiard system becomes superintegrable when the mass
ratio takes the form $m/M = \tan^2(\pi/q)$ for integer
$q$~\cite{Davis2017}.  At these special values, the system possesses a
third independent integral---the Chevalley invariant---beyond the
energy and angular momentum.

Table~\ref{tab:masses} shows that the dimension sequence is invariant
across all tested mass ratios, including both superintegrable values
($q = 3, 4, 5, 6$) and a generic non-superintegrable configuration.
In every case, $d(k) = [3, 6, 17, 116]$.

\begin{table}[h]
\caption{1D Calogero-Moser dimension sequences at
Galperin superintegrable mass ratios and generic masses.}
\label{tab:masses}
\begin{ruledtabular}
\begin{tabular}{lcccccc}
$m/M$ & $q$ & SI? & $d(0)$ & $d(1)$ & $d(2)$ & $d(3)$ \\
\colrule
3       & 3 & Yes & 3 & 6 & 17 & 116 \\
1$^\ddagger$  & 4 & Yes & 3 & 6 & 17 & 116 \\
$\approx 0.53$ & 5 & Yes & 3 & 6 & 17 & 116 \\
$1/3$   & 6 & Yes & 3 & 6 & 17 & 116 \\
$(1, 2.7, 0.4)$ & --- & No & 3 & 6 & 17 & 116 \\
\end{tabular}
\end{ruledtabular}
\footnotetext{$^\ddagger$Equal masses; coincides with the base
computation in Table~\ref{tab:gaps}.}
\end{table}


% ==================================================================
\section{Discussion}
\label{sec:discussion}
% ==================================================================

\subsection{The Singularity Dichotomy}

Our results establish a clear classification principle: the Poisson
algebra dimension sequence of a three-body system with pairwise
interactions depends on the \emph{singularity class} of the pair
potential, not on its integrability properties.

\begin{itemize}
\item \textbf{Singular potentials} ($V \sim r^{-n}$, $\log r$,
  Yukawa): The algebra is infinite-dimensional with the universal
  growth sequence $d(k) = [3, 6, 17, 116, \ldots]$.  This holds
  regardless of the exponent $n$, spatial dimension $d$, particle
  masses, or coupling constants.

\item \textbf{Regular potentials} ($V \sim r^2$): The algebra is
  finite-dimensional, saturating at $d = 15$ by level 3.
\end{itemize}

This dichotomy is robust: the transition occurs at the boundary
between potentials that are singular at $r = 0$ and those that are
analytic there.  For power-law singularities $V \propto r^{-n}$,
the substitution $u = 1/r$ yields potentials polynomial in~$u$, but
the universality extends to transcendental cases: $\log r = -\log u$
and the Yukawa potential $e^{-\mu r}/r = u\,e^{-\mu/u}$ are not
polynomial in~$u$, yet produce the same dimension sequence.  What
all singular potentials share is the chain rule
(\ref{eq:chainrule})---the cubic relation $\partial u/\partial q
\propto u^3$ is determined entirely by the definition $u = 1/r$ and
is independent of the specific form of $V(u)$.  The harmonic
potential $V = r^2$ requires no auxiliary variables and generates a
finite algebra because its Poisson brackets close after finitely many
iterations.  Since the chain rule is shared, the dimension sequence
is determined by the combinatorial structure of the pairwise
interaction graph---how many bodies interact and in which
pairs---rather than by the specific potential acting along each edge.

\subsection{Why Integrability is Invisible}

The Calogero-Moser model possesses two functionally independent
conserved quantities beyond the total energy: the squared angular
momentum $L^2$ (ensuring Liouville integrability) and, at
superintegrable mass ratios, the Chevalley invariant $J$ (ensuring
maximal superintegrability).  Why does the pairwise Poisson algebra
fail to ``see'' these integrals?

The answer lies in the nature of these conserved quantities.  The
angular momentum $L^2 = (M+m)^2(Yw - yW)^2$ is a \emph{global}
function of the full phase space---it does not decompose into pairwise
contributions.  The pairwise Hamiltonians $H_{ij}$ generate an algebra
that respects the combinatorial structure of the interaction graph
(which pairs of bodies interact), but does not capture collective
conservation laws that emerge from the geometry of the full
configuration space.

In the language of algebraic geometry, the pairwise Poisson algebra
is generated by a specific \emph{subvariety} of the full algebra of
observables, one that is determined by the pair structure of the
Hamiltonian.  The integrals of motion of the Calogero model lie
outside this subvariety---they are ``invisible'' to the pairwise
generators, even though they constrain the dynamics.

This explains the universality.  The chain rule
table~(\ref{eq:chainrule}) is determined entirely by the definition
$u_{ij} = 1/r_{ij}$ and is identical for every potential---Newtonian,
Calogero, or otherwise.  All Poisson brackets among the $H_{ij}$
therefore share the same derivative structure; what differs is only
the form of the potential term $V(u)$ in the initial generators.
For regular potentials such as $V = r^2$, no auxiliary variables
enter and the brackets close after finitely many iterations.  For
any potential singular at $r = 0$, the $u$-substitution introduces
the universal cubic chain rule, and iterated brackets produce new
functionally independent expressions at each level.  The dimension
count reflects the combinatorial structure of the pairwise
interaction graph filtered through this shared chain rule, not the
specific potential or any global dynamical property such as
integrability.

\subsection{Connection to Galperin Billiards}

Galperin's billiard system~\cite{Galperin2003} computes $\pi$ through
a counting argument that relies on the conservation of $L^2$ in the
equivalent Calogero model.  Davis \textit{et al.}~\cite{Davis2017}
showed that this conservation law makes the system integrable and,
at special mass ratios, superintegrable.  Our computation confirms
that these powerful conservation laws leave no imprint on the pairwise
Poisson algebra: the billiard system generates the same algebraic
structure as generic, non-integrable gravity.

This result is not a deficiency of our method but rather a precise
characterization of what the pairwise algebra measures.  It is an
\emph{algebraic} invariant of the singularity class, not a
\emph{dynamical} invariant of the trajectory structure.


% ==================================================================
\section{Conclusion}
\label{sec:conclusion}
% ==================================================================

We have computed the Poisson algebra dimension sequence for the
one-dimensional Calogero-Moser three-body system---an exactly
integrable model motivated by Galperin's billiard construction---and
found $d(k) = [3, 6, 17, 116]$, identical to every other singular
pair potential tested.  This result holds at all Galperin
superintegrable mass ratios ($q = 3, 4, 5, 6$) and for generic mass
configurations.

Together with the contrasting behavior of the harmonic oscillator
($d(k) \to 15$), this establishes that the Poisson algebra dimension
sequence is a \emph{singularity class invariant}: it distinguishes
potentials singular at the origin from regular ones, but is blind to
integrability, superintegrability, mass ratios, spatial dimension,
and coupling constants.

The universal sequence $[3, 6, 17, 116, \ldots]$ for singular
potentials suggests a deeper algebraic structure---perhaps related to
the combinatorics of the pairwise interaction graph---that warrants
further investigation.  In particular, the sequence's relationship to
known integer sequences and its behavior at higher filtration levels
($k \geq 4$, where computational costs become substantial) remain
open questions.  We note that the cross-potential universality
established here holds through level~3; whether the sequences for
different singular potentials continue to agree at level~4 and beyond,
or eventually diverge, is the most important open computational test
of the singularity-class hypothesis.

\begin{acknowledgments}
Computational resources were provided by Amazon Web Services (EC2).
Symbolic computations used SymPy; numerical linear algebra used
NumPy and SciPy.
\end{acknowledgments}


\begin{thebibliography}{10}

\bibitem{Galperin2003}
G.~Galperin,
\textit{Playing pool with $\pi$ (the number $\pi$ from a billiard
point of view)},
Regul.\ Chaotic Dyn.\ \textbf{8}, 375 (2003).

\bibitem{Davis2017}
B.~Davis, D.~Geyer, and A.~Ravishanker,
\textit{The dynamics of digits: Calculating pi with Galperin's
billiards},
arXiv:1712.06698 (2017).

\bibitem{Calogero1971}
F.~Calogero,
\textit{Solution of the one-dimensional $N$-body problems with
quadratic and/or inversely quadratic pair potentials},
J.\ Math.\ Phys.\ \textbf{12}, 419 (1971).

\bibitem{Moser1975}
J.~Moser,
\textit{Three integrable Hamiltonian systems connected with isospectral
deformations},
Adv.\ Math.\ \textbf{16}, 197 (1975).

\bibitem{Calogero1969}
F.~Calogero,
\textit{Solution of a three-body problem in one dimension},
J.\ Math.\ Phys.\ \textbf{10}, 2191 (1969).

\bibitem{Sutherland1972}
B.~Sutherland,
\textit{Exact results for a quantum many-body problem in one
dimension},
Phys.\ Rev.\ A \textbf{4}, 2019 (1971);
\textit{ibid.}\ \textbf{5}, 1372 (1972).

\bibitem{Olshanetsky1983}
M.~A.~Olshanetsky and A.~M.~Perelomov,
\textit{Quantum integrable systems related to Lie algebras},
Phys.\ Rep.\ \textbf{94}, 313 (1983).

\bibitem{Wojciechowski1983}
S.~Wojciechowski,
\textit{Superintegrability of the Calogero-Moser system},
Phys.\ Lett.\ A \textbf{95}, 279 (1983).

\bibitem{paper1}
[Authors],
\textit{The Poisson algebra of the three-body problem:
dimension sequence, mass invariance, and potential comparison},
(2026, in preparation).

\bibitem{paper3}
[Authors],
\textit{Universal dimension sequences of pairwise Poisson algebras:
independence from spatial dimension, potential exponent, and charge
sign},
(2026, in preparation).

\end{thebibliography}

\end{document}
